% GENERATED FILE. Edit canonical lessons, then run npm run generate:books.
% canonical-source-hash: fnv1a64:1e68fe99f027cc2a
% canonical-lessons: SA-C126-vistrtah, SA-C126-sankirnah, SA-C126-gabhirah, SA-C126-ardrah, SA-C126-suskah

\chapter{Wide, Narrow, Deep, Wet, Dry}
\label{ch:sa-wide-narrow-deep-wet-dry}
\hlchaptermodality{126}

\begin{chapteropening}
\textbf{By the end of this chapter:} \emph{I can say wide, narrow, deep, wet and dry.}

Complete the last lesson of chapter 126: i can say wide, narrow, deep, wet and dry.
\end{chapteropening}

\section[\texorpdfstring{\sk{विस्तृतः} (vistṛtaḥ)}{विस्तृतः (vistṛtaḥ)}]{\sk{विस्तृतः} (vistṛtaḥ) — wide}

\label{lesson:SA-C126-vistrtah}



\begin{quote}
Before the new one: say the Sanskrit for clean, then the Sanskrit for dirty.
\end{quote}

\subsection*{You'll want to know: \sk{विस्तृतः}}
\textbf{\sk{विस्तृतः}} — \emph{vistṛtaḥ} — \textquotedblleft{}wide\textquotedblright{}.

\begin{grammarlens}[title={What you've built}]
One more word for this chapter.
\end{grammarlens}

\subsection*{Your turn}
\begin{itemize}
\raggedright
  \item \emph{Say it:} \emph{vistṛtaḥ}
  \item \emph{Say it:} \emph{vistṛtaḥ}, once more
  \item \emph{Say it:} say \emph{vistṛtaḥ}
  \item \emph{Recall:} say the Sanskrit for clean, then the Sanskrit for dirty, then say \emph{vistṛtaḥ} again
\end{itemize}

\begin{tcolorbox}[breakable,colback=teal!4,colframe=teal!35!black,title={Before you move on}]
What does \sk{विस्तृतः} mean? (\textquotedblleft{}Wide\textquotedblright{}.) Say it once more.
\end{tcolorbox}
\section[\texorpdfstring{\sk{संकीर्णः} (saṅkīrṇaḥ)}{संकीर्णः (saṅkīrṇaḥ)}]{\sk{संकीर्णः} (saṅkīrṇaḥ) — narrow}

\label{lesson:SA-C126-sankirnah}



\begin{quote}
Before the new one: say the Sanskrit for dirty, then the Sanskrit for wide.
\end{quote}

\subsection*{You'll want to know: \sk{संकीर्णः}}
\textbf{\sk{संकीर्णः}} — \emph{saṅkīrṇaḥ} — \textquotedblleft{}narrow\textquotedblright{}.

\begin{grammarlens}[title={What you've built}]
One more word for this chapter.
\end{grammarlens}

\subsection*{Your turn}
\begin{itemize}
\raggedright
  \item \emph{Say it:} \emph{saṅkīrṇaḥ}
  \item \emph{Say it:} \emph{saṅkīrṇaḥ}, once more
  \item \emph{Say it:} say \emph{saṅkīrṇaḥ}
  \item \emph{Recall:} say the Sanskrit for dirty, then the Sanskrit for wide, then say \emph{saṅkīrṇaḥ} again
\end{itemize}

\begin{tcolorbox}[breakable,colback=teal!4,colframe=teal!35!black,title={Before you move on}]
What does \sk{संकीर्णः} mean? (\textquotedblleft{}Narrow\textquotedblright{}.) Say it once more.
\end{tcolorbox}
\section[\texorpdfstring{\sk{गभीरः} (gabhīraḥ)}{गभीरः (gabhīraḥ)}]{\sk{गभीरः} (gabhīraḥ) — deep}

\label{lesson:SA-C126-gabhirah}



\begin{quote}
Before the new one: say the Sanskrit for wide, then the Sanskrit for narrow.
\end{quote}

\subsection*{You'll want to know: \sk{गभीरः}}
\textbf{\sk{गभीरः}} — \emph{gabhīraḥ} — \textquotedblleft{}deep\textquotedblright{}.

\begin{grammarlens}[title={What you've built}]
One more word for this chapter.
\end{grammarlens}

\subsection*{Your turn}
\begin{itemize}
\raggedright
  \item \emph{Say it:} \emph{gabhīraḥ}
  \item \emph{Say it:} \emph{gabhīraḥ}, once more
  \item \emph{Say it:} say \emph{gabhīraḥ}
  \item \emph{Recall:} say the Sanskrit for wide, then the Sanskrit for narrow, then say \emph{gabhīraḥ} again
\end{itemize}

\begin{tcolorbox}[breakable,colback=teal!4,colframe=teal!35!black,title={Before you move on}]
What does \sk{गभीरः} mean? (\textquotedblleft{}Deep\textquotedblright{}.) Say it once more.
\end{tcolorbox}
\section[\texorpdfstring{\sk{आर्द्रः} (ārdraḥ)}{आर्द्रः (ārdraḥ)}]{\sk{आर्द्रः} (ārdraḥ) — wet}

\label{lesson:SA-C126-ardrah}



\begin{quote}
Before the new one: say the Sanskrit for narrow, then the Sanskrit for deep.
\end{quote}

\subsection*{You'll want to know: \sk{आर्द्रः}}
\textbf{\sk{आर्द्रः}} — \emph{ārdraḥ} — \textquotedblleft{}wet\textquotedblright{}.

\begin{grammarlens}[title={What you've built}]
One more word for this chapter.
\end{grammarlens}

\subsection*{Your turn}
\begin{itemize}
\raggedright
  \item \emph{Say it:} \emph{ārdraḥ}
  \item \emph{Say it:} \emph{ārdraḥ}, once more
  \item \emph{Say it:} say \emph{ārdraḥ}
  \item \emph{Recall:} say the Sanskrit for narrow, then the Sanskrit for deep, then say \emph{ārdraḥ} again
\end{itemize}

\begin{tcolorbox}[breakable,colback=teal!4,colframe=teal!35!black,title={Before you move on}]
What does \sk{आर्द्रः} mean? (\textquotedblleft{}Wet\textquotedblright{}.) Say it once more.
\end{tcolorbox}
\section[\texorpdfstring{\sk{शुष्कः} (śuṣkaḥ)}{शुष्कः (śuṣkaḥ)}]{\sk{शुष्कः} (śuṣkaḥ) — dry}

\label{lesson:SA-C126-suskah}



\begin{quote}
Before the new one: say the Sanskrit for deep, then the Sanskrit for wet.
\end{quote}

\subsection*{You'll want to know: \sk{शुष्कः}}
\textbf{\sk{शुष्कः}} — \emph{śuṣkaḥ} — \textquotedblleft{}dry\textquotedblright{}.

\begin{grammarlens}[title={What you've built}]
One more word for this chapter.
\end{grammarlens}

\subsection*{Your turn}
\begin{itemize}
\raggedright
  \item \emph{Say it:} \emph{śuṣkaḥ}
  \item \emph{Say it:} \emph{śuṣkaḥ}, once more
  \item \emph{Say it:} say \emph{śuṣkaḥ}
  \item \emph{Recall:} say the Sanskrit for deep, then the Sanskrit for wet, then say \emph{śuṣkaḥ} again
\end{itemize}

\begin{tcolorbox}[breakable,colback=teal!4,colframe=teal!35!black,title={Before you move on}]
What does \sk{शुष्कः} mean? (\textquotedblleft{}Dry\textquotedblright{}.) Say it once more.
\end{tcolorbox}
